\chapter{Q as Value Compression, Policy as Preference Compression}
\label{ch:compression}

Most treatments of RL describe the Q function and the policy as two different
``things to learn.''  This chapter offers a more illuminating framing: both are
\textbf{compressors}, but they compress different objects.  Understanding what
each one compresses---and how the compressed representations connect---explains
why algorithms are structured the way they are and why the field moved from
value-first methods to policy-first methods.


%% ================================================================
\section{Q: Compressing Sparse Rewards into Dense Value}
\label{sec:q-compression}

Consider a typical RL environment.  The agent receives a reward of $+1$ only
when it reaches a distant goal; every other timestep yields $r=0$.  Raw rewards
are \emph{sparse}---they carry almost no per-step information.

The Q function compresses the entire future into a single scalar:
\begin{equation}
  Q(s,a) = \mathbb{E}\!\left[\sum_{k=0}^{\infty}\gamma^k\, r_{t+k}
  \;\Big|\; s_t = s,\; a_t = a\right].
\end{equation}

This is a lossy compression.  The full distribution of possible futures---all
the branching trajectories, stochastic transitions, and delayed rewards---is
collapsed into one number.  But it is an extraordinarily useful compression:
it tells you, at every state--action pair, how much long-term value to expect.

\begin{keyidea}[Q as Value Compressor]
The Q function compresses sparse rewards and complex future rollouts into a
dense, per-step scalar signal.  It transforms a credit-assignment problem
(``which past action caused this delayed reward?'') into a local lookup
(``what is this action worth right here?'').
\end{keyidea}


%% ================================================================
\section{Policy: Compressing Behavioral Preferences}
\label{sec:policy-compression}

The policy is also a compressor, but its target is different.  Rather than
compressing \emph{value}, it compresses \emph{preference}---the agent's
behavioral tendency in each state.

A stochastic policy $\pi(a|s)$ represents preferences as a probability
distribution.  A deterministic policy $\mu(s)$ represents preferences as a
single point.  Both are compressed summaries of ``how I prefer to act here,''
distilled from experience.

\begin{keyidea}[Policy as Preference Compressor]
The policy compresses the agent's accumulated behavioral preferences into a
compact representation---a distribution $\pi(a|s)$ or a point $\mu(s)$---that
can be evaluated in a single forward pass.
\end{keyidea}


%% ================================================================
\section{Two Compression Strategies}
\label{sec:two-strategies}

The distinction between value compression and preference compression maps
directly onto the two major families of RL algorithms.

\subsection{Value-First: Compress Value, Then Extract Preference}

Q-learning and DQN follow a two-stage strategy:
\begin{enumerate}[nosep]
  \item \textbf{Compress value}: learn $Q(s,a)$ for all actions.
  \item \textbf{Extract preference}: apply $\operatorname{argmax}$ to convert
    the value landscape into a behavioral choice.
\end{enumerate}

\begin{equation}
  \underbrace{\text{Sparse rewards} \;\longrightarrow\; Q(s,a)}_{\text{value compression}}
  \;\longrightarrow\;
  \underbrace{a^{\ast} = \operatorname{argmax}_a Q(s,a)}_{\text{preference extraction}}
\end{equation}

The advantage is that $Q$ can be learned from off-policy data with high sample
efficiency.  The disadvantage is that the second stage---argmax---is a fragile
bridge between value and preference (Section~\ref{sec:fragile-argmax}).

\subsection{Policy-First: Compress Preference Directly}

Policy gradient methods skip the value-compression stage (or use it only as an
auxiliary signal) and directly optimize the preference representation:
\begin{equation}
  \underbrace{\text{Sampled returns / advantages}
  \;\longrightarrow\; \pi_\theta(a|s)}_{\text{preference compression}}
\end{equation}

The advantage is robustness: there is no brittle argmax step.  The disadvantage
is data hunger---on-policy methods must discard data whenever the policy
changes, and the gradient signal from sampled returns is noisy.

\subsection{Actor-Critic: Value Compressor Guides Preference Compressor}

Actor-Critic architectures maintain \emph{both} compressors simultaneously:
\begin{itemize}[nosep]
  \item The \textbf{critic} compresses value (or advantage).
  \item The \textbf{actor} compresses preference.
  \item The critic's compressed value signal \emph{guides} the actor's
    preference updates.
\end{itemize}

\begin{equation}
  \underbrace{Q_\phi(s,a) \text{ or } A_\phi(s,a)}_{\text{value compressor (critic)}}
  \;\xrightarrow{\text{guides}}\;
  \underbrace{\pi_\theta(a|s)}_{\text{preference compressor (actor)}}
\end{equation}

This is not simply ``one network scores, another network acts.''  It is a
cooperative compression scheme in which the value compressor provides a
low-variance training signal that accelerates the preference compressor's
learning.

The risk is mutual contamination: if the critic's value compression is poor
(overestimation, dead regions), it corrupts the actor's preference compression.
This is why Actor-Critic methods need careful stabilization---target networks,
double critics, entropy bonuses---to keep both compressors healthy.


%% ================================================================
\section{Why Hard Argmax Is a Fragile Bridge}
\label{sec:fragile-argmax}

The argmax operation is the standard way to convert a Q-value landscape into a
decision:
\begin{equation}
  a^{\ast} = \operatorname{argmax}_a Q(s,a).
\end{equation}

In theory, if $Q = Q^{\ast}$, this recovers the optimal policy.  In practice,
$Q$ is an approximation with noise, and argmax has two pathological properties:

\begin{enumerate}
  \item \textbf{Discontinuity.}  Argmax is a discontinuous function of $Q$.  A
    tiny perturbation in Q-values can flip the selected action entirely.  This
    means that small estimation errors produce large behavioral changes.
  \item \textbf{Noise amplification.}  Among $k$ noisy estimates, argmax
    preferentially selects the one with the largest positive noise.  This is the
    overestimation bias discussed in Chapter~\ref{ch:bias-management}, but seen
    here as a \emph{compression failure}: the value compressor's noise is
    amplified at the value-to-preference interface.
\end{enumerate}

A soft policy $\pi(a|s) \propto \exp(Q(s,a)/\alpha)$ replaces the hard argmax
with a smooth, temperature-controlled mapping.  Small errors in $Q$ produce
small changes in the distribution, not discontinuous action switches.  This is
one of the deepest reasons the field moved toward stochastic, soft policies.

\begin{keyidea}[Argmax as Compression Bottleneck]
Hard argmax is a fragile bridge between value compression and preference
compression.  It amplifies noise and is discontinuous.  Soft policies replace
this brittle interface with a smooth, differentiable mapping---making the
overall compression pipeline more robust.
\end{keyidea}


%% ================================================================
\section{Algorithm Comparison Through the Compression Lens}
\label{sec:compression-table}

Table~\ref{tab:compression} summarizes major algorithms in terms of what they
compress, how they make decisions, and what can go wrong.

\begin{table}[ht]
\centering
\caption{Algorithms viewed as compression systems.}
\label{tab:compression}
\begin{tabular}{@{}llll@{}}
\toprule
\textbf{Algorithm} & \textbf{Compression object} &
\textbf{Decision method} & \textbf{Core risk} \\
\midrule
Q-learning / DQN &
  Long-term value $Q$ &
  Hard argmax &
  Q noise amplified by argmax \\
Double DQN / TD3 &
  Value (stabilized) &
  Conservative argmax &
  Underestimation; slower learning \\
Policy Gradient &
  Preference $\pi$ directly &
  Sample from $\pi$ &
  High variance; data hungry \\
PPO &
  Preference $\pi$ (stabilized) &
  Clipped update of $\pi$ &
  May be overly conservative \\
SAC &
  Soft preference &
  Sample from entropy-regularized $\pi$ &
  Objective rewritten by entropy \\
Actor-Critic &
  $Q/A$ (value) + $\pi$ (preference) &
  Critic guides actor &
  Compressors can contaminate each other \\
\bottomrule
\end{tabular}
\end{table}


%% ================================================================
\section{The Compression Perspective in One Sentence}
\label{sec:compression-summary}

\begin{keyidea}[The Compression View of RL]
$Q$ is the compression of sparse rewards and long-term value into a dense
scalar.  $\pi$ is the compression of behavioral preferences into a
distribution.  The history of RL algorithm design is the story of learning how
to build, connect, and stabilize these two compressors---and of discovering
that directly compressing preferences is often more robust than compressing
value first and then extracting preferences through the fragile bottleneck of
argmax.
\end{keyidea}
